
\documentclass[10pt,aspectratio=169]{beamer}
\usepackage[utf8]{inputenc}
\usepackage[T1]{fontenc}
\usepackage[catalan]{babel}
\usepackage{amsmath,amssymb,booktabs,array,multicol}
\usepackage{tikz,pgfplots,xcolor,listings}
\usetikzlibrary{positioning,arrows.meta,calc,shapes.geometric}
\pgfplotsset{compat=1.18}
\usepgfplotslibrary{fillbetween}

\definecolor{UABGreen}{RGB}{0,103,71}
\definecolor{UABLight}{RGB}{224,241,236}
\definecolor{DeepBlue}{RGB}{20,65,110}
\definecolor{AccentOrange}{RGB}{230,126,34}
\definecolor{SoftGray}{RGB}{245,247,250}
\definecolor{CodeBack}{RGB}{248,248,248}
\definecolor{CodeFrame}{RGB}{210,215,220}

\usetheme{Madrid}
\usecolortheme[named=UABGreen]{structure}
\setbeamercolor{title}{fg=white,bg=UABGreen}
\setbeamercolor{frametitle}{fg=white,bg=UABGreen}
\setbeamercolor{block title}{fg=white,bg=DeepBlue}
\setbeamercolor{block body}{fg=black,bg=SoftGray}
\setbeamercolor{alerted text}{fg=AccentOrange}
\setbeamertemplate{navigation symbols}{}
\setbeamertemplate{itemize item}{\color{UABGreen}$\blacktriangleright$}
\setbeamertemplate{itemize subitem}{\color{AccentOrange}$\bullet$}
\setbeamertemplate{enumerate item}{\color{UABGreen}\insertenumlabel.}
\setbeamertemplate{footline}{%
  \leavevmode%
  \hbox{%
  \begin{beamercolorbox}[wd=.72\paperwidth,ht=2.4ex,dp=1ex,leftskip=1.4em]{author in head/foot}%
    Estadística -- Grau en Enginyeria Informàtica
  \end{beamercolorbox}%
  \begin{beamercolorbox}[wd=.28\paperwidth,ht=2.4ex,dp=1ex,center]{date in head/foot}%
    \insertframenumber{} / \inserttotalframenumber
  \end{beamercolorbox}}%
  \vskip0pt%
}

\lstdefinelanguage{R}{%
  morekeywords={if,else,repeat,while,function,for,in,next,break,TRUE,FALSE,NULL,NA,NaN,Inf,
    library,mean,median,sd,var,summary,table,prop.table,xtabs,addmargins,round,
    ggplot,aes,geom_point,geom_line,geom_col,geom_histogram,geom_density,geom_smooth,
    geom_tile,facet_wrap,labs,theme_minimal,summarise,group_by,mutate,read_csv,filter,
    arrange,set.seed,rbinom,rgeom,rpois,runif,rexp,rnorm,dnorm,pnorm,qnorm,dt,pt,qt,
    dchisq,pchisq,qchisq,df,pf,qf,dbinom,pbinom,dpois,ppois,replicate,sample,rep,seq,
    tibble,data.frame,confint,t.test,var.test,chisq.test,binom.test,prop.test,lm,
    broom,glance,augment,tidy,future_map_dbl,plan,multisession,map_dbl,purrr,future,furrr},
  sensitive=true,
  morecomment=[l]\#,
  morestring=[b]",
  morestring=[b]'
}
\lstset{language=R,basicstyle=\ttfamily\scriptsize,keywordstyle=\color{DeepBlue}\bfseries,
  commentstyle=\color{gray},stringstyle=\color{AccentOrange},backgroundcolor=\color{CodeBack},
  frame=single,rulecolor=\color{CodeFrame},breaklines=true,columns=fullflexible,keepspaces=true,
  showstringspaces=false,
  literate={à}{{\`a}}1 {è}{{\`e}}1 {é}{{\'e}}1 {í}{{\'i}}1 {ï}{{\"i}}1 {ò}{{\`o}}1 {ó}{{\'o}}1 {ú}{{\'u}}1 {ü}{{\"u}}1 {ç}{{\c{c}}}1 {À}{{\`A}}1 {È}{{\`E}}1 {É}{{\'E}}1 {Í}{{\'I}}1 {Ò}{{\`O}}1 {Ó}{{\'O}}1 {Ú}{{\'U}}1}

\newcommand{\R}{\textsf{R}}
\newcommand{\RStudio}{\textsf{RStudio/Posit}}
\newcommand{\GitHub}{\textsf{GitHub}}
\newcommand{\important}[1]{\textcolor{UABGreen}{\textbf{#1}}}
\newcommand{\exemple}[1]{\textcolor{AccentOrange}{\textbf{#1}}}
\newcommand{\Prob}{\mathbb{P}}
\newcommand{\E}{\mathbb{E}}
\newcommand{\Var}{\operatorname{Var}}
\newcommand{\IC}{\operatorname{IC}}

\title[Inferència VI]{Inferència VI}
\author{Estadística -- Grau en Enginyeria Informàtica}
\date{Curs 2026--27}

\begin{document}
\begin{frame}[plain]
  \titlepage
  \vspace{-0.45cm}
  \begin{center}
  \small\textcolor{DeepBlue}{Estadística computacional amb \R, simulació i exemples d'enginyeria informàtica}
  \end{center}
\end{frame}


\begin{frame}[fragile]{Per què tests no paramètrics?}
\begin{itemize}
  \item Fins ara hem contrastat paràmetres: mitjanes, proporcions i variàncies.
  \item Ara contrastem estructures: independència o ajust a una distribució.
  \item Són especialment útils amb variables categòriques i taules de contingència.
\end{itemize}
\begin{block}{Exemples informàtics}
Sistema operatiu i tipus d'error; navegador i conversió; distribució esperada de codis d'estat HTTP.
\end{block}
\end{frame}

\begin{frame}[fragile]{Objectius}
\begin{enumerate}
  \item Formular tests $\chi^2$ d'independència i d'ajustament.
  \item Calcular valors observats i esperats.
  \item Interpretar residus i contribucions.
  \item Fer els tests amb \texttt{chisq.test()} en \R.
  \item Relacionar resultats amb decisions sobre producte o sistema.
\end{enumerate}
\end{frame}

\begin{frame}[fragile]{Dos tipus de tests $\chi^2$}
\begin{columns}[T]
\column{0.48\textwidth}
\begin{block}{Independència}
\[
H_0:\text{dues variables són independents.}
\]
Exemple: tipus de dispositiu i error.
\end{block}
\column{0.48\textwidth}
\begin{block}{Ajustament}
\[
H_0:\text{les dades segueixen una distribució donada.}
\]
Exemple: freqüències de codis HTTP segons proporcions esperades.
\end{block}
\end{columns}
\vspace{0.3cm}
En ambdós casos, l'estadístic és gran quan observat i esperat difereixen molt.
\end{frame}

\begin{frame}[fragile]{Test d'independència}
\begin{block}{Taula de contingència}
Per cada cel·la $(i,j)$ tenim observat $o_{ij}$ i esperat
\[
e_{ij}=\frac{(\text{total fila }i)(\text{total columna }j)}{n}.
\]
\[
E=\sum_{i,j}\frac{(o_{ij}-e_{ij})^2}{e_{ij}}
\sim \chi^2_{(f-1)(c-1)}.
\]
\end{block}
\begin{alertblock}{Condició pràctica}
Els esperats haurien de ser prou grans, habitualment $e_{ij}\ge 5$.
\end{alertblock}
\end{frame}

\begin{frame}[fragile]{Exemple: dispositiu i error}
\begin{columns}[T]
\column{0.45\textwidth}
\begin{tabular}{lrr}
\toprule
 & Error & Sense error\\
\midrule
Mòbil & 48 & 952\\
Escriptori & 22 & 978\\
\bottomrule
\end{tabular}
\vspace{0.3cm}
\[
H_0:\text{dispositiu i error són independents.}
\]
\column{0.48\textwidth}
\begin{lstlisting}
taula <- matrix(c(48, 952, 22, 978),
                nrow = 2, byrow = TRUE)
rownames(taula) <- c("mobil", "desktop")
colnames(taula) <- c("error", "ok")
chisq.test(taula, correct = FALSE)
\end{lstlisting}
\end{columns}
\end{frame}

\begin{frame}[fragile]{Mirar més enllà del p-valor}
\begin{lstlisting}
res <- chisq.test(taula, correct = FALSE)
res$expected      # valors esperats
res$residuals     # residus de Pearson
res$stdres        # residus estandarditzats
\end{lstlisting}
\begin{itemize}
  \item Les cel·les amb residus grans expliquen la dependència.
  \item Això ajuda a traduir el test en accions concretes.
\end{itemize}
\begin{block}{Interpretació}
No només interessa si hi ha relació, sinó on apareix la diferència.
\end{block}
\end{frame}

\begin{frame}[fragile]{Test d'ajustament}
\begin{block}{Una sola variable categòrica}
Tenim freqüències observades $o_i$ i probabilitats esperades $p_i$.
\[
e_i=n p_i,
\qquad
E=\sum_i\frac{(o_i-e_i)^2}{e_i}.
\]
\end{block}
\[
H_0:\text{les dades segueixen les probabilitats }p_i.
\]
\begin{lstlisting}
observat <- c(200, 720, 80)
esperat_p <- c(.2, .75, .05)
chisq.test(observat, p = esperat_p)
\end{lstlisting}
\end{frame}

\begin{frame}[fragile]{Exemple: codis HTTP}
\begin{columns}[T]
\column{0.50\textwidth}
En condicions normals esperem:
\begin{itemize}
  \item 2xx: 92\%,
  \item 3xx: 5\%,
  \item 4xx: 2\%,
  \item 5xx: 1\%.
\end{itemize}
Avui observem 10.000 respostes: 9000, 520, 260, 220.
\column{0.44\textwidth}
\begin{lstlisting}
obs <- c(9000, 520, 260, 220)
p0 <- c(.92, .05, .02, .01)
chisq.test(obs, p = p0)
\end{lstlisting}
\begin{alertblock}{Pregunta}
El patró d'avui és compatible amb el funcionament normal?
\end{alertblock}
\end{columns}
\end{frame}

\begin{frame}[fragile]{Visualització de contribucions}
\begin{lstlisting}
res <- chisq.test(obs, p = p0)
contrib <- (obs - res$expected)^2 / res$expected

tibble(codi = c("2xx","3xx","4xx","5xx"), contrib) |>
  ggplot(aes(codi, contrib)) +
  geom_col() +
  theme_minimal()
\end{lstlisting}
\begin{block}{Idea}
Les contribucions indiquen quines categories expliquen el rebuig de $H_0$.
\end{block}
\end{frame}

\begin{frame}[fragile]{Limitacions}
\begin{itemize}
  \item Requereix freqüències esperades prou grans.
  \item No mesura causalitat.
  \item Amb mostres molt grans, diferències petites poden ser significatives.
  \item Cal complementar amb mida de l'efecte i interpretació de domini.
\end{itemize}
\begin{block}{Opcional}
Per taules $2\times 2$, es pot mirar també odds ratio o diferència de proporcions.
\end{block}
\end{frame}

\end{document}
